\chapter{Improve Your Sleep}

\section*{Definition and Overview}
Sleep is essential for physical and mental health, and poor sleep affects mood, concentration, immunity, and long-term health. Improving sleep usually involves a combination of habits known as sleep hygiene, an environment that supports rest, and attention to the daytime and evening routines that influence sleep. For most people with short-term sleep problems, simple changes make a real difference. When sleep problems persist for weeks or months, they may indicate insomnia or another sleep disorder that benefits from structured treatment, including cognitive behavioural therapy for insomnia, which is the most effective long-term treatment.

\section*{Epidemiology}
Sleep problems are extremely common. Around one in three adults experiences short-term insomnia, and up to one in ten has chronic insomnia. Poor sleep is more common with age, stress, shift work, and chronic illness, and it contributes to a large burden of accidents, reduced work performance, and health problems including heart disease, obesity, and depression. Most people who report poor sleep have never discussed it with a doctor or received structured treatment, despite effective options being available. Public health campaigns increasingly recognise sleep as a fundamental pillar of health alongside diet and exercise.

\section*{Aetiology and Pathophysiology}
Sleep is controlled by the interaction of two systems: the circadian rhythm, an internal 24-hour clock that regulates the timing of sleepiness and wakefulness, and the homeostatic drive, which builds sleep pressure the longer a person is awake. Sleep quality and quantity are influenced by light exposure, activity, food and caffeine, stress, and the sleep environment. Insomnia typically develops when a trigger, such as stress or illness, disrupts sleep, and is then maintained by habits such as worrying in bed, irregular sleep schedules, and daytime napping that train the brain to associate bed with wakefulness. These mechanisms explain why behavioural treatment, rather than sleeping tablets alone, is the cornerstone of lasting improvement.

\section*{Clinical Features}
\begin{itemize}
    \item Difficulty falling asleep, staying asleep, or waking too early
    \item Waking feeling unrefreshed, despite enough time in bed
    \item Daytime effects: fatigue, irritability, poor concentration, and low mood
    \item Heavy reliance on coffee, energy drinks, or daytime naps to get through the day
    \item Sleeping later on weekends to catch up
    \item Lying in bed for long periods worrying about not sleeping
    \item Physical signs of sleep deprivation: heavy eyelids, yawning, reduced motivation, and lowered immunity
\end{itemize}

\section*{Differential Diagnosis}
Persistent poor sleep may be due to conditions other than simple insomnia, which need different treatment.
\begin{itemize}
    \item \textbf{Obstructive sleep apnoea:} loud snoring with pauses in breathing, gasping, and daytime sleepiness
    \item \textbf{Depression and anxiety:} commonly cause early waking and difficulty getting to sleep
    \item \textbf{Restless legs syndrome:} an urge to move the legs that interferes with sleep onset
    \item \textbf{Chronic pain, reflux, and medical illness:} any condition that disturbs sleep
    \item \textbf{Medicines and substances:} caffeine, alcohol, stimulants, and some medicines disrupt sleep
    \item \textbf{Shift work and jet lag:} circadian rhythm disorders from an irregular schedule
\end{itemize}

\section*{When to Seek Medical Care}
See a doctor if poor sleep is persistent, affecting daytime function, or accompanied by loud snoring, gasping, or pauses in breathing, or by persistent low mood. Seek review before using over-the-counter sleep aids regularly, and ask about structured insomnia treatment. Sleep problems in children that are severe or associated with snoring, breathing difficulty, or behavioural problems also warrant assessment.

\textbf{Call 000 (or local emergency services) for:}
\begin{itemize}
    \item Sudden severe shortness of breath at night or while lying flat, which may indicate heart or lung problems
    \item A baby or child who is very difficult to wake or is breathing abnormally during sleep
    \item Confusion, collapse, or a seizure
\end{itemize}

\section*{Diagnosis and Investigations}
\begin{itemize}
    \item \textbf{Sleep history:} sleep schedule, patterns, environment, and daytime effects, often supported by a sleep diary
    \item \textbf{Screening questionnaires:} tools such as the Insomnia Severity Index and Epworth Sleepiness Scale
    \item \textbf{Review of medicines, caffeine, alcohol, and stress:} important contributors to poor sleep
    \item \textbf{Sleep study (polysomnography):} overnight monitoring in a laboratory or at home where sleep apnoea is suspected
    \item \textbf{Actigraphy:} a wrist monitor that measures sleep-wake patterns over days to weeks
    \item \textbf{Assessment of underlying conditions:} blood tests and review for depression, thyroid disease, or pain
\end{itemize}

\section*{Management}
\begin{itemize}
    \item \textbf{Sleep hygiene:} consistent sleep and wake times, a cool dark quiet bedroom, and avoiding screens and stimulants before bed
    \item \textbf{Cognitive behavioural therapy for insomnia (CBT-I):} the first-line treatment, combining stimulus control, sleep restriction, and restructuring of unhelpful thoughts about sleep
    \item \textbf{Stimulus control:} using the bed only for sleep, and getting up if unable to sleep
    \item \textbf{Sleep restriction:} limiting time in bed to match actual sleep, then gradually increasing it
    \item \textbf{Relaxation techniques:} progressive muscle relaxation, mindfulness, and breathing exercises
    \item \textbf{Short-term medicines:} occasionally used for a few weeks, but not recommended for long-term insomnia
    \item \textbf{Treatment of underlying conditions:} sleep apnoea, pain, reflux, and mood disorders
    \item \textbf{Light therapy and schedule adjustment:} for shift workers and circadian problems
\end{itemize}

\section*{Complications}
\begin{itemize}
    \item Daytime fatigue, poor concentration, and reduced work or study performance
    \item Increased risk of road and workplace accidents
    \item Mood disturbance, anxiety, and depression
    \item Impaired immunity and increased susceptibility to infection
    \item Longer-term associations with obesity, diabetes, high blood pressure, and heart disease
    \item Dependence on sleeping tablets with long-term use
\end{itemize}

\section*{Prognosis and Outcomes}
Short-term sleep problems usually resolve as the trigger passes and habits improve. Chronic insomnia is treatable: CBT-I improves sleep in most people and has lasting effects, unlike sleeping tablets, which lose effectiveness over time and carry risks. Around two-thirds of people with chronic insomnia respond well to structured treatment. Even when a sleep disorder such as apnoea is present, effective treatment dramatically improves sleep quality and reduces the associated health risks.

\section*{Prevention and Self-Care}
\begin{itemize}
    \item Keep a regular sleep-wake schedule, including on weekends
    \item Get morning light exposure and be physically active during the day
    \item Avoid caffeine after midday and alcohol close to bedtime
    \item Wind down with a relaxing routine in the hour before bed
    \item Keep the bedroom cool, dark, and quiet, and use it only for sleep
    \item Put screens away at least 30--60 minutes before bed
    \item Get up if you cannot sleep and return to bed when drowsy
    \item Avoid long or late daytime naps
    \item Seek professional help early if problems persist beyond a few weeks
\end{itemize}

\section*{Sources and Further Reading}
The following reputable sources provide further information for healthcare students and patients seeking reliable, current advice about improving sleep.
\begin{itemize}
    \item Sleep Health Foundation. \textit{Tips for better sleep and insomnia information.} https://www.sleephealthfoundation.org.au/
    \item Healthdirect Australia. \textit{Insomnia and sleep problems.} https://www.healthdirect.gov.au/
    \item Beyond Blue. \textit{Sleep and mental health.} https://www.beyondblue.org.au/
    \item NHS. \textit{Insomnia and sleep hygiene.} https://www.nhs.uk/conditions/insomnia/
    \item Mayo Clinic. \textit{Insomnia treatment and self-care.} https://www.mayoclinic.org/
    \item Sleep Foundation. \textit{Sleep hygiene and healthy sleep.} https://www.sleepfoundation.org/
\end{itemize}
